%% kapa: exam=1 round=36 year=2025 period=2 form=A subject=accounting subject_ko=회계학 seq=23 number=63
%% source: Q-Net 감정평가사 기출문제 (cst00309) artlSeq=5231525 file=raw/1008/5231525/2025년 제36회 감정평가사 1차시험 2교시 문제지 원본.pdf
%% conversion: kordoc 4.24.51 -> tools/gen1.py (auto); review=proofread-sonnet; flags=none
\begin{question}{63}
(주)감평의 20×1년 보통주에 귀속되는 당기순이익은 ￦1,000,000, 가중평균유통보통주식수는 100주, 중단사업손익은 없다. (주)감평이 희석주당이익을 계산할 때 필요한 자료는 다음과 같다.
\par{\small\centering
\begin{tabular}{|c|c|c|}\hline
종류 & \parbox[c]{0.3\linewidth}{\centering\vspace{2pt}보통주로 전환(행사)되었다고 가정할 경우 추가적으로 유통되었을 가중평균유통보통주식수 (분모)\vspace{2pt}} & \parbox[c]{0.3\linewidth}{\centering\vspace{2pt}보통주로 전환(행사)되었다고 가정할 경우 증가하는 법인세 효과를 반영한 보통주귀속이익 (분자)\vspace{2pt}} \\ \hline
전환사채 & 150주 & ￦425,000 \\ \hline
전환우선주 & 100 & 190,000 \\ \hline
신주인수권부사채 & 200 & 10,000 \\ \hline
옵션 & 100 & 0 \\ \hline
\end{tabular}
\par}
(주)감평의 20×1년 희석주당이익은? (단, 20×1년 잠재적보통주식의 변동은 없다.)
\begin{choices}
  \item ￦2,400
  \item ￦2,500
  \item ￦2,525
  \item ￦5,000
  \item ￦10,000
\end{choices}
\end{question}
\answer{1}
